\documentclass[11pt]{article}
\usepackage[margin=1in]{geometry}
\usepackage{amsmath,amssymb,amsthm}
\usepackage[hidelinks]{hyperref}

\newtheorem{theorem}{Theorem}
\newtheorem{lemma}[theorem]{Lemma}

\title{Odd-Size Antipodal Rational 5-Designs for the Chebyshev Measure: Resolution of Mishima--Lu--Sawa--Uchida Problem~8.3}
\author{Salem Levinstein}
\date{November 2026}

\begin{document}
\maketitle

\begin{abstract}
We resolve Problem~8.3 of Mishima, Lu, Sawa, and Uchida
(\texttt{arXiv:2503.21151v4}; \emph{Res.\ Number Theory} 12 (2026), No.~39) by
exhibiting an explicit $37$-point antipodal rational $5$-design for the
Chebyshev measure $(1-t^2)^{-1/2}\,dt/\pi$ on $(-1,1)$ with common denominator
$240$. Together with explicit base designs of sizes $35$ and $39$ and disjoint
$6$-point rational blocks, we prove that an odd-size antipodal rational
$5$-design for the Chebyshev measure exists if and only if $n \ge 35$,
completing the existence spectrum of rational antipodal $5$-designs. The entire
result, including a self-contained $2$-adic proof of non-existence for
$n \le 33$, is formally verified in Lean~4.
\end{abstract}

\section{Introduction}

An $n$-point \emph{antipodal rational $5$-design} for the Chebyshev measure
$d\mu(t) = (1 - t^2)^{-1/2}\,dt/\pi$ on $(-1, 1)$ is a finite subset
$X = -X \subset (-1, 1) \cap \mathbb{Q}$ of cardinality $n$ such that
\[
\frac{1}{n} \sum_{x \in X} x^k = \int_{-1}^{1} t^k\,d\mu(t)
\qquad (0 \le k \le 5).
\]
Odd moments vanish automatically by antipodality, while the second and fourth
Chebyshev moments are $1/2$ and $3/8$. When $n = 2N + 1$ is odd, writing
$X = \{0\} \cup \{\pm x_1, \dots, \pm x_N\}$ with distinct
$x_i \in (0, 1) \cap \mathbb{Q}$, the design condition is equivalent to the
Diophantine system
\begin{equation}\label{eq:hk}
\sum_{i=1}^{N} x_i^2 = \frac{2N + 1}{4},
\qquad
\sum_{i=1}^{N} x_i^4 = \frac{3(2N + 1)}{16}.
\end{equation}

In~\cite{MLSU2026}, Mishima, Lu, Sawa, and Uchida determined the complete
spectrum of even sizes $2N$, proved that no odd-size design exists for
$2N + 1 \le 33$ (Theorem~8.1), constructed odd-size designs for
$N \equiv 5 \pmod 6$ with $N \ge 17$ (Theorem~8.2), and posed:
\begin{quote}
\emph{\textbf{Problem 8.3.} Find an antipodal $5$-design with $37$ rational
points, as the smallest open case.}
\end{quote}
We settle Problem~8.3 and complete the odd-size classification.

\begin{theorem}\label{thm:main}
An antipodal rational $5$-design for the Chebyshev measure with an odd number of
points $n = 2N + 1$ exists if and only if $n \ge 35$ ($N \ge 17$). In
particular, a $37$-point design exists:
\begin{align*}
X_{37} = \{0\} \cup \Bigl\{\pm \frac{a}{240} : a \in \{&15,57,61,69,93,129,141,147,151,\\
&183,199,209,211,215,234,235,238,239\}\Bigr\}.
\end{align*}
\end{theorem}

\section{Existence for Every $N \ge 17$}

We first exhibit three base positive half-sets $S_{N_0} \subset (0, 1) \cap \mathbb{Q}$
solving~\eqref{eq:hk} for $N_0 \in \{17, 18, 19\}$:
\begin{itemize}
\item $N_0 = 17$ ($n = 35$, $D = 1092$): numerators\\
$\{55, 87, 171, 305, 533, 575, 733, 753, 765, 881, 905, 966, 981, 985, 1083, 1085, 1089\}$;
\item $N_0 = 18$ ($n = 37$, $D = 240$, Problem~8.3): numerators\\
$\{15, 57, 61, 69, 93, 129, 141, 147, 151, 183, 199, 209, 211, 215, 234, 235, 238, 239\}$;
\item $N_0 = 19$ ($n = 39$, $D = 420$): numerators\\
$\{47, 93, 108, 113, 129, 153, 161, 305, 309, 343, 353, 357, 365, 366, 369, 395, 405, 408, 413\}$.
\end{itemize}
Direct integer arithmetic confirms $\sum a_i^2 = \frac{2N_0+1}{4} D^2$ and
$\sum a_i^4 = \frac{3(2N_0+1)}{16} D^4$ in all three cases.

\begin{lemma}[Monotone disjoint $6$-point blocks]\label{lem:blocks}
For $t \in \mathbb{N}$ with $t \ge 20$, set
\[
a(t) = \frac{2t^2 - 22t - 13}{14(t^2 + t + 1)},
\quad
b(t) = \frac{13t^2 + 4t - 11}{14(t^2 + t + 1)},
\quad
c(t) = \frac{11t^2 + 26t + 2}{14(t^2 + t + 1)}.
\]
Then $a(t)^2 + b(t)^2 + c(t)^2 = 3/2$, $a(t)^4 + b(t)^4 + c(t)^4 = 9/8$, and
\[
0 < a(t) < \tfrac{2}{14} < \tfrac{11}{14} < c(t) < \tfrac{12}{14} < b(t) < \tfrac{13}{14} < 1.
\]
Moreover, $a(t)$ and $b(t)$ are strictly increasing and $c(t)$ is strictly
decreasing for $t \ge 20$.
\end{lemma}

\begin{proof}
Clearing the common denominator $14(t^2 + t + 1)$ reduces both moment equations
and the six interval inequalities to elementary polynomial identities and
quadratic inequalities valid for $t \ge 20$. For $20 \le t_1 < t_2$,
\begin{align*}
a(t_2) - a(t_1) &= \frac{(t_2 - t_1)(24 t_1 t_2 + 15(t_1 + t_2) - 9)}{14(t_1^2+t_1+1)(t_2^2+t_2+1)} > 0, \\
b(t_2) - b(t_1) &= \frac{(t_2 - t_1)(9 t_1 t_2 + 24(t_1 + t_2) + 15)}{14(t_1^2+t_1+1)(t_2^2+t_2+1)} > 0, \\
c(t_1) - c(t_2) &= \frac{(t_2 - t_1)(15 t_1 t_2 - 9(t_1 + t_2) - 24)}{14(t_1^2+t_1+1)(t_2^2+t_2+1)} > 0. \qedhere
\end{align*}
\end{proof}

\begin{proof}[Proof of Existence ($N \ge 17$)]
Given any finite solution $S \subset (0, 1) \cap \mathbb{Q}$ of~\eqref{eq:hk},
the strict monotonicity in Lemma~\ref{lem:blocks} implies that each of
$a, b, c$ takes values in $S$ for at most $|S|$ integers $t \ge 20$. Hence among
the $3|S| + 1$ integers $20 \le t \le 20 + 3|S|$, at least one $t$ yields a
triple $\{a(t), b(t), c(t)\}$ disjoint from $S$. Adjoining it increases $|S|$ by
$3$, $\sum x^2$ by $3/2$, and $\sum x^4$ by $9/8$, preserving~\eqref{eq:hk}.
Inducting from $S_{17}$, $S_{18}$, or $S_{19}$ according to $N \bmod 3$ covers
every $N \ge 17$.
\end{proof}

\section{Non-Existence for $N \le 16$}

For completeness we give a short $2$-adic proof of Theorem~8.1
of~\cite{MLSU2026}. Suppose $S = \{x_1, \dots, x_N\}$ satisfies~\eqref{eq:hk}
with $N \le 16$ and $M = 2N + 1$ odd. Choose the minimal common denominator
$D \ge 1$ with $x_i = y_i / D$ ($y_i \in \mathbb{Z}$) so that if $2 \mid D$ then
some $y_i$ is odd. Then
\[
4\sum_{i=1}^N y_i^2 = M D^2, \qquad 16\sum_{i=1}^N y_i^4 = 3 M D^4.
\]
Let $r = \#\{i : y_i \text{ odd}\}$; then $\sum y_i^4 \equiv r \pmod{16}$.
\begin{itemize}
\item If $2 \nmid D$, then $M D^2$ is odd, contradicting $4 \mid M D^2$.
\item If $D = 2 D_1$ with $D_1$ odd, then $\sum y_i^2 = M D_1^2 \equiv M \pmod 8$
and $\sum y_i^4 = 3 M D_1^4 \equiv 3M \pmod 8$. Letting $j = \#\{i : y_i \equiv 2 \pmod 4\}$,
we obtain $r \equiv 3M \pmod 8$ and $r + 4j \equiv M \pmod 8$, whence
$2M + 4j \equiv 0 \pmod 8$, contradicting $M$ odd.
\item If $D = 4 D_2$, then $\sum y_i^4 = 48 M D_2^4 \equiv 0 \pmod{16}$, so
$16 \mid r$. Since $1 \le r \le N \le 16$, we have $r = N = 16$ ($M = 33$), and
every $y_i$ is odd. Writing $y_i^2 = 1 + 8 t_i$, $T_1 = \sum t_i$,
$T_2 = \sum t_i^2$ yields $16 + 8 T_1 = 4 M D_2^2$ and
$16 + 16 T_1 + 64 T_2 = 48 M D_2^4$. If $D_2$ is odd, the first equation fails
modulo $8$ ($0 \not\equiv 4$). If $D_2$ is even, the first equation gives
$16 \mid 16 + 8 T_1$ ($T_1$ even), while the second gives
$256 \mid 16(1 + T_1 + 4 T_2)$ ($T_1$ odd), a contradiction. \qed
\end{itemize}

\section{Formal and Computational Verification}

Theorem~\ref{thm:main} (both Problem~8.3 and the full odd-size biconditional
\texttt{mishima\_odd\_\allowbreak 5design\_iff}) is formally verified in Lean~4
(module \texttt{Conjectures/\allowbreak MishimaLuSawaUchida5Design.lean}) with
zero \texttt{sorry} and standard axioms only. An independent Python script
(\texttt{verify.py}) verifies all base designs, block identities, explicit
constructions for $17 \le N \le 80$, and the $2$-adic residue checks.

\begin{thebibliography}{9}
\bibitem{MLSU2026}
T.~Mishima, X.-N.~Lu, M.~Sawa, and Y.~Uchida,
\emph{Geometric designs and Hilbert--Kamke equations of degree five for
classical orthogonal polynomials},
Res.\ Number Theory \textbf{12} (2026), Paper No.~39;
\texttt{arXiv:2503.21151v4}.
\end{thebibliography}

\end{document}
